\documentclass[11pt]{article}
\usepackage{icmlko}
\begin{document}

\title{모델은 이미 방향을 배웠다 --- 축이 서른여섯이면 가릴 이유가
한쪽에서 사라진다}
\author{월드모델 실험실}
\date{노트 465 · 466 · 2026년 8월 2일}
\maketitle

\begin{abstract}
노트 463 이 \texttt{OOC\_DROP} 을 미결로 내리며 남긴 물음은 두 개였다 ---
\emph{두 짝의 갈림이 진짜인가}, 그리고 \emph{진짜라면 무엇이 부호를
정하나}.

\textbf{하나 --- 진짜다.} 각 짝을 무작위 세 겹으로 갈랐다. KR 겹 셋이
\emph{전부} 음수($-0.0567 \cdot -0.0697 \cdot -0.0490$)고 앱 겹 셋이
\emph{전부} 양수($+0.0349 \cdot +0.0407 \cdot +0.0507$)다. 수집 안 폭
$0.0207$ 대 수집 사이 간격 $\mathbf{0.1005}$ --- \textbf{$4.9$배}.
2짝 계기는 잡음을 보고 있지 않다.

\textbf{둘 --- 축 수는 아니다.} 라벨을 안 보고 중복도 큰 축부터 하나씩
더 빼 봤다. 부호가 만나기는커녕 \textbf{각자 자기 방향으로 더 커진다}
--- KR $-0.0584 \to -0.0955$, 앱 $+0.0435 \to +0.0830$. \textbf{같은 축
수($4$)에서 KR 은 $-0.0584$, 앱은 $+0.0830$} 이다. 다른 축들은 효과를
\emph{희석}하고 있었을 뿐이다.

\textbf{셋 --- 모델의 쓰임이 정한다.}

\begin{center}
\begin{tabular}{lrrrr}
\hline
짝 & 짝 안 축$\leftrightarrow$라벨 & \textbf{모델 쓰임$\leftrightarrow$축값} & 모델 쓰임$\leftrightarrow$라벨 & 가림 효과 \\
\hline
KR 만화 & $-0.320$ & $\mathbf{-0.604}$ & $\mathbf{+0.183}$ & $-0.0584$ \\
비게임 앱 & $-0.566$ & $\mathbf{-0.026}$ & $-0.152$ & $+0.0435$ \\
\hline
\end{tabular}
\end{center}

\textbf{36축 모델은 KR 에서 그 축을 이미 음의 방향으로 쓴다}($-0.604$) ---
KR 자신의 부호($-0.320$)와 맞다. 그러니 모델의 쓰임이 라벨과 \emph{양으로}
붙고($+0.183$) \textbf{가리면 손해다.}
\end{abstract}

\begin{figure}[t]
\centering
\includegraphics[width=\textwidth]{figs/f1.pdf}
\caption{\textbf{가} --- 겹 여섯이 자기 수집 부호를 지킨다. \textbf{나} ---
축을 빼도 부호가 안 만난다. \textbf{다} --- 짝 안 부호는 둘 다 음수라
못 가르는데, \emph{모델의 쓰임}은 갈린다.}
\end{figure}

\section{노트 441 의 전제가 축 다섯짜리였다}

노트 441 은 이렇게 적었다 --- \emph{``모델은 큰 양수 쪽(모바일 $+0.648$)을
배워 처음 보는 수집에 거꾸로 쓴다.''} 그래서 가리는 게 옳았다.

\textbf{36축에서는 그 전제가 깨진다.} 모델이 KR 에서 그 축을 쓰는 방향이
$-0.604$ --- \emph{음수}다. 깊이 $12$ 와 서른여섯 축이 있으면 나무는 그
축을 다른 축들과 묶어 \textbf{조건부로} 쓰고, KR 처럼 생긴 행에서는
제대로 음의 방향으로 쓴다. \emph{가릴 이유가 사라진 것이다.}

앱에서는 왜 안 그런가. 모델의 쓰임이 축값에 대해 \textbf{거의 평평하다}
($-0.026$) --- 값에 따라 단조로 움직이지 않고 상호작용이 끌고 간다. 그
움직임이 앱의 라벨과 어긋나므로($-0.152$) \textbf{가리는 게 이득이다.}

\begin{quote}
\textbf{읽는 법} --- 가림은 \emph{축이 나쁠 때}가 아니라 \textbf{모델이
그 축을 그 수집에서 못 쓸 때} 이롭다. 축은 KR 과 앱 모두에서 음의 부호를
갖는데(둘 다 $-0.3 \sim -0.6$), 모델은 KR 에서만 그 방향을 맞힌다.
\end{quote}

\section{라벨 없이 고를 수 있는 길이 보인다}

\texttt{OOC\_DROP} 이 미결인 이유는 \emph{영샷}이기 때문이다 --- 새 수집의
라벨이 없으니 짝 안 부호를 못 잰다. 그런데 \textbf{``모델 쓰임$\leftrightarrow$
축값''은 라벨이 필요 없다.} 예보 둘(축 켬 · 끔)과 축값만 있으면 된다.

두 짝에서 그 값이 $-0.604$ 와 $-0.026$ 으로 크게 갈리고, 부호도 정확히
갈린다. \textbf{후보 규칙} --- \emph{모델이 그 축을 축값에 단조롭게
쓰면(절댓값이 크면) 그대로 두고, 거의 안 쓰면(0 근처면) 가린다.}

\textbf{짝이 둘이라 이 규칙은 아직 가설이다.} 미리 적어 둔다 --- 다음에
집 밖 짝이 하나 더 생기면 \textbf{라벨을 보기 전에} 이 값으로 예보하고,
맞는지 본다.

\section*{남은 것}

\begin{itemize}
\item \textbf{규칙의 문턱} --- $-0.604$ 와 $-0.026$ 사이 어디인지 두
점으로는 못 정한다. 노트 466 의 축 빼기 팔들(짝마다 셋)로 안쪽 변이를
만들 수 있나.
\item \textbf{셋째 짝} --- CN 만화는 \texttt{entry\_friction} 이 아예
없어 이 물음에 표를 못 던진다. 태거 갈래다.
\item \textbf{다른 축에도 같은 물음} --- 지금은 \texttt{entry\_friction}
하나만 본다. 다른 넷도 수집마다 모델의 쓰임이 갈리나.
\end{itemize}

\end{document}
